\documentclass[11pt]{amsart}
\usepackage{mathmemo}
\renewcommand{\memotitle}{OPAC-009 solution}
\renewcommand{\memodate}{September 9, 2026}

%\usepackage[T1]{fontenc}
%\usepackage{lmodern}
%\usepackage{amsmath,amssymb,amsthm,mathtools}
%\usepackage[margin=1.05in]{geometry}
%\usepackage{microtype}
%\usepackage[colorlinks=true,linkcolor=blue!50!black,citecolor=blue!50!black,urlcolor=blue!50!black]{hyperref}
\usepackage{xcolor}
%\newtheorem{theorem}{Theorem}[section]
%\newtheorem{proposition}[theorem]{Proposition}
%\newtheorem{lemma}[theorem]{Lemma}
%\newtheorem{corollary}[theorem]{Corollary}
%\theoremstyle{remark}
%\newtheorem{remark}[theorem]{Remark}
%\newcommand{\R}{\mathbb R}
\renewcommand{\C}{\mathcal C}
\newcommand{\SP}{\mathcal{SP}}
\newcommand{\dom}{\trianglerighteq}
\newcommand{\wt}{\operatorname{wt}}
\newcommand{\tops}{\tau}
\newcommand{\bots}{\beta}
\newcommand{\un}{\uplus}
\newcommand{\coef}[2]{[\delta^{#1}]\,#2}

\title[Extreme rays of the cone of log-concavity]{Extreme rays of the cone of log-concavity}

\author{Claude (Fable-5.1), with ChatGPT-6 (Astra)}
 \thanks{Directed by D.~E.~Anderson, \texttt{anderson.2804@math.osu.edu}}

\date{September 9, 2026}
\begin{document}
\begin{abstract}
White conjectured that a product of Schur functions indexed by partitions
with at most two parts spans an extreme ray of the cone of log-concavity
if and only if its factors form a nested multiset. We prove the conjecture.
The main ingredient is a specialization of the complete homogeneous
symmetric functions to a positive log-concave sequence. Making its
logarithmic increments constant on an interval, and then perturbing them,
turns a positive decomposition of Schur products into an identity between
products of linear forms. Interval counts constrain the endpoints of
these forms. A suitable coefficient then singles out the original nested
multiset, with dominance supplying the final comparison when a one-part
factor is present.
\end{abstract}
\maketitle

\section{Introduction}

Let $\Lambda=\R[h_1,h_2,\ldots]$ be the ring of symmetric functions,
with $\deg h_i=i$ and $h_0=1$. For $N\ge1$, let $\SP_N$ be the set of
multisets of nonempty partitions with at most two parts and total size $N$.
For $A\in\SP_N$, write
\[
 s_A=\prod_{\rho\in A}s_\rho,
 \qquad
 \C_N=\operatorname{cone}\{s_A:A\in\SP_N\}\subseteq\Lambda_N.
\]
Jacobi--Trudi gives
\begin{equation}\label{eq:jt}
 s_{(a,b)}=h_ah_b-h_{a+1}h_{b-1}\quad(a\ge b\ge1),
 \qquad s_{(c)}=h_c.
\end{equation}
Thus $\C_N$ is the degree-$N$ part of White's \emph{cone of
log-concavity}. Its generators are nonnegative on every positive
log-concave sequence $(h_i)_{i\ge0}$. The question is which of these
generators are necessary: which products span extreme rays?

We call a two-part member $(a,b)$ of $A$ a \emph{pair}, with \emph{top}
$a$ and \emph{bottom} $b$, and a one-part member $(c)$ a
\emph{singleton}. Following White, two members \emph{interlace} if,
after possibly exchanging them, they are one of the following:
\begin{enumerate}
\item pairs $(a,b)$ and $(d,e)$ with $a>d\ge b>e$;
\item a pair $(a,b)$ and a singleton $(c)$ with $a>b$ and $a\ge c\ge b$;
\item two singletons.
\end{enumerate}
A multiset is \emph{nested} if no two of its members interlace.
Repeated members are counted with multiplicity. In particular, a pair
$(p,p)$ interlaces with nothing, and a nested multiset has at most one
singleton.

\begin{theorem}\label{thm:main}
For every $N\ge1$ and $A\in\SP_N$, the product $s_A$ spans an extreme
ray of $\C_N$ if and only if $A$ is nested.
\end{theorem}

This proves White's Conjecture~3, listed as OPAC-009 \cite{OPAC}. White proved the
necessity of nestedness and established the converse when all the parts
of the members of $A$ are distinct \cite[Theorems~4 and~12]{White}.
Our proof allows arbitrary repetitions.

Here is the idea. A specialization $h_k\mapsto g(k)>0$ with constant
logarithmic increments on $[x,y+1]$ kills exactly the factors
$s_{(a,b)}$ with $x\le b\le a\le y$. A perturbation of these increments
replaces each vanishing factor, to first order, by a positive constant
times $\delta_b-\delta_{a+1}$. Positivity prevents cancellation of the
lowest powers of the perturbation parameter. Consequently, any positive
decomposition of $s_A$ forces each competitor to have at least as many
pairs as $A$ in every interval.

These interval inequalities control the sorted lists of tops and
bottoms. If $A$ has no singleton, the coefficient of its all-bottom
monomial finishes the proof. If $A$ has a singleton $(c)$, nestedness
splits its pairs into blocks below, at, and above $c$. We select the
top endpoints below $c$ and the bottom endpoints above $c$; dominance
then forces both blocks, and the singleton, to agree with $A$.

\section{The reduction to nested competitors}

For a multiset $A$, let $\phi(A)$ be the partition obtained by sorting
all its parts. Write $\mu\dom\lambda$ for dominance order, with zeros
appended when needed. White's Schur-support statement is
\begin{equation}\label{eq:support}
 s_A=\sum_{\mu\dom\phi(A)}c_A^\mu s_\mu,
 \qquad c_A^\mu\ge0,\qquad c_A^{\phi(A)}=1
 \quad\text{\cite[Proposition~6]{White}.}
\end{equation}
We use it together with White's theorem that non-nested products are
not extreme. No result from the withdrawn paper \cite{GMTWYZ} is needed.

\begin{lemma}\label{lem:distinct}
Distinct multisets in $\SP_N$ give nonproportional products $s_A$.
\end{lemma}
\begin{proof}
The polynomial $h_ah_b-h_{a+1}h_{b-1}$ is primitive and linear in
$h_{a+1}$ over the polynomial ring in the other variables: its two
coefficients are coprime. Gauss's lemma therefore makes it irreducible.
The polynomials $h_c$ are also irreducible. These irreducibles are
pairwise nonassociate: a two-part factor has largest variable index
$a+1$, and its coefficient of that variable is $-h_{b-1}$ (or $-1$ if
$b=1$), which determines $b$ and its normalization. A one-part factor
cannot be associate to a two-term polynomial. Unique factorization now
proves the assertion. All these arguments take place in a polynomial
ring with finitely many variables for any given products.
\end{proof}

\begin{lemma}\label{lem:reduction}
If $A$ is nested and $s_A$ does not span an extreme ray, then
\begin{equation}\label{eq:decomp}
 s_A=\sum_{B\in\mathcal S}c_Bs_B,
 \qquad c_B>0,
\end{equation}
where $\mathcal S$ is a nonempty finite set of nested multisets
$B\ne A$, each satisfying $\phi(B)\dom\phi(A)$.
\end{lemma}
\begin{proof}
The cone $\C_N$ is finitely generated and pointed, since it lies in the
Schur-positive cone. It is therefore generated by its extreme rays,
each of which contains a generator $s_B$. Such a $B$ is nested by
White's theorem and Lemma~\ref{lem:distinct}. Express $s_A$ as a
positive combination of these generators. None is proportional to
$s_A$, since $s_A$ is assumed not to span an extreme ray.
Finally, the coefficient of $s_{\phi(B)}$ on the right is positive.
There is no cancellation in the Schur basis, so
\eqref{eq:support} forces $\phi(B)\dom\phi(A)$.
\end{proof}

\section{An interval specialization}

For an integer interval $I=[x,y]$, with $1\le x\le y$, let $B_I$
be the multiset of pairs $(a,b)\in B$ with $x\le b\le a\le y$, and set
\[
 n_I(B)=|B_I|,
 \qquad
 P_{B_I}(\delta)=\prod_{(a,b)\in B_I}(\delta_b-\delta_{a+1}).
\]
The variables of this polynomial are indexed by $T=[x,y+1]$.
An empty product is $1$.

\begin{lemma}[Interval degeneration]\label{lem:interval}
Suppose \eqref{eq:decomp} holds, without requiring $A$ or the $B$'s
to be nested. Fix $I=[x,y]$ and put $m=n_I(A)$. Then
\begin{equation}\label{eq:count}
 n_I(B)\ge m\qquad(B\in\mathcal S).
\end{equation}
Moreover, there are constants $K_A,K_B>0$, independent of $\delta$,
such that
\begin{equation}\label{eq:level}
 K_AP_{A_I}(\delta)=
 \sum_{\substack{B\in\mathcal S\\n_I(B)=m}}
 c_BK_BP_{B_I}(\delta).
\end{equation}
\end{lemma}

\begin{proof}
Choose a nonincreasing sequence $d^0(k)$, $k\ge1$, whose only equal
consecutive values occur within $T$: it is constant on $T$ and strictly
decreases at every other step. For $e=(e_k)_{k\in T}$, extended by zero
outside $T$, set
\[
 d^t(k)=d^0(k)+te_k,
 \qquad g_t(k)=\exp\left(\sum_{j=1}^k d^t(j)\right),
 \qquad \psi_{t,e}(h_k)=g_t(k).
\]
Here $g_t(0)=1$, and algebraic independence of the $h_k$ makes
$\psi_{t,e}$ a well-defined ring homomorphism. By \eqref{eq:jt},
\begin{equation}\label{eq:factor}
 \psi_{t,e}(s_{(a,b)})
 =g_t(a)g_t(b)\left(1-\exp(d^t(a+1)-d^t(b))\right).
\end{equation}
At $t=0$ this is zero exactly when $(a,b)$ lies inside $I$; otherwise
it is positive. For a pair inside $I$, Taylor expansion gives
\[
 \psi_{t,e}(s_{(a,b)})
 =t\,g_0(a)g_0(b)(e_b-e_{a+1})+O(t^2).
\]
Multiplying the factors, for each fixed $e$ we obtain
\begin{equation}\label{eq:leading}
 \psi_{t,e}(s_B)
 =t^{n_I(B)}K_BP_{B_I}(e)+O(t^{n_I(B)+1}),
\end{equation}
where
\[
 K_B=
 \prod_{(a,b)\in B_I}g_0(a)g_0(b)
 \prod_{\substack{(a,b)\in B\\(a,b)\notin B_I}}
       \psi_{0,e}(s_{(a,b)})
 \prod_{(c)\in B}g_0(c)>0.
\]
The value at $t=0$ is independent of $e$, so $K_B$ is too. Formula
\eqref{eq:leading} remains valid when some differences $e_b-e_{a+1}$
vanish; it does not assert an exact order of vanishing in those directions.

Let $r=\min_{B\in\mathcal S}n_I(B)$. If $r<m$, choose $e_k$
strictly decreasing on $T$. Every $P_{B_I}(e)$ is then positive.
Apply $\psi_{t,e}$ to \eqref{eq:decomp}, divide by $t^r$, and let
$t\to0$. The left side tends to zero, whereas the right side tends to
$\sum_{n_I(B)=r}c_BK_BP_{B_I}(e)>0$. This proves \eqref{eq:count}.

Now take arbitrary $e$, divide the specialized identity by $t^m$,
and let $t\to0$. We get \eqref{eq:level} evaluated at $e$.
Since this holds for every $e\in\R^T$, it is a polynomial identity.
\end{proof}

The count inequality holds for \emph{every} interval in one and the same
decomposition. The constants in \eqref{eq:level} may change with the
interval; no comparison between constants from different intervals is
required.

\begin{corollary}[Insertion of a rectangular pair]\label{cor:rectangle}
If $s_D$ spans an extreme ray, then so does $s_{D\un\{(p,p)\}}$
for every $p\ge1$.
\end{corollary}
\begin{proof}
In a positive decomposition of $s_Ds_{(p,p)}$, apply
\eqref{eq:count} with $I=[p,p]$. Every term contains $(p,p)$.
Cancel one copy of its Schur factor in the polynomial ring.
A decomposition using other multisets would give a decomposition of
$s_D$ using other multisets, contrary to extremality and
Lemma~\ref{lem:distinct}.
\end{proof}

\section{Endpoints and their monomials}

For a multiset $M$ of $m$ pairs, let $\tops(M)$ and $\bots(M)$ be its
multisets of tops and bottoms. List them increasingly as
$a_1(M)\le\cdots\le a_m(M)$ and $b_1(M)\le\cdots\le b_m(M)$.
Write $\wt(M)=\sum_i a_i(M)+\sum_i b_i(M)$ and
\[
 P_M=\prod_{(a,b)\in M}(\delta_b-\delta_{a+1}),
 \qquad \delta^X=\prod_{v\in X}\delta_v.
\]
In $\delta^{\tops(M)+1}$, the $1$ is added to every member of the
multiset $\tops(M)$.

\begin{lemma}[Endpoint inequalities]\label{lem:endpoints}
Let $D,M$ consist of $m$ pairs inside $[w,u]$. If
$n_J(M)\ge n_J(D)$ for every subinterval $J\subseteq[w,u]$, then
\begin{equation}\label{eq:endpoints}
 b_i(M)\ge b_i(D),\qquad a_i(M)\le a_i(D)
 \qquad(1\le i\le m).
\end{equation}
\end{lemma}
\begin{proof}
The intervals $[x,u]$ count bottoms at least $x$, and the intervals
$[w,y]$ count tops at most $y$. Their inequalities are precisely the two
coordinatewise comparisons between the increasing lists.
\end{proof}

\begin{lemma}[Endpoint monomials]\label{lem:monomials}
Let $D,M$ consist of $m$ pairs.
\begin{enumerate}
\item If $\sum\bots(M)\ge\sum\bots(D)$, then
\[
 \coef{\bots(D)}{P_M}=
 \begin{cases}1&\text{if }\bots(M)=\bots(D),\\0&\text{otherwise.}\end{cases}
\]
\item If $\sum\tops(M)\le\sum\tops(D)$, then
\[
 \coef{\tops(D)+1}{P_M}=
 \begin{cases}(-1)^m&\text{if }\tops(M)=\tops(D),\\0&\text{otherwise.}\end{cases}
\]
\end{enumerate}
\end{lemma}
\begin{proof}
Give $\delta_k$ weight $k$. In expanding $P_M$, choosing
$\delta_{a+1}$ instead of $\delta_b$ strictly increases the weight,
since $a+1>b$. Thus the unique choice of least weight selects every
bottom; its monomial is $\delta^{\bots(M)}$, with coefficient $1$.
Under the first hypothesis, the target has weight no greater than this
minimum, so it can occur only as that term. The unique choice of
greatest weight selects every top plus one and has coefficient
$(-1)^m$. The second assertion follows in the same way. Repeated
factors cause no difficulty, since every change of choice changes the
weight strictly.
\end{proof}

\begin{lemma}[Reconstruction]\label{lem:reconstruction}
A nested multiset of pairs is determined by its tops and its bottoms.
\end{lemma}
\begin{proof}
Let $b$ be the largest bottom, with multiplicity $r$. We claim that the
$r$ tops paired with $b$ are the $r$ smallest available tops at least
$b$. Indeed, let $a^*$ be the largest top paired with $b$. If a pair
$(a,d)$ has $d<b\le a<a^*$, then $(a^*,b)$ and $(a,d)$ interlace.
Thus every top at least $b$ that is paired with a smaller bottom is
at least $a^*$. This proves the claim, including ties.
Remove the $r$ pairs just determined, and repeat. Induction reconstructs
the entire multiset.
\end{proof}

\section{Proof of the main theorem}

White's theorem proves the non-nested direction. Suppose $A$ is nested
and $s_A$ is not extreme. Lemma~\ref{lem:reduction} gives a
decomposition \eqref{eq:decomp} with nested $B\ne A$ and
$\phi(B)\dom\lambda:=\phi(A)$.

Let $D$ be the multiset of pairs in $A$, put $m=|D|$, and let
$n=\ell(\lambda)$. Nestedness gives $n\in\{2m,2m+1\}$.
Let $w$ and $u$ be the smallest and largest parts of $\lambda$,
and put $I=[w,u]$. Dominance implies $\ell(\phi(B))\le n$:
a partition with more than $n$ positive parts cannot dominate a
partition of the same size with $n$ parts. Hence the interval lemma gives
\begin{equation}\label{eq:length}
 2m\le2n_I(B)\le\ell(\phi(B))\le n\le2m+1.
\end{equation}
It follows that $B$ has exactly $m$ pairs, all inside $I$, and at most
one singleton. Denote its pair multiset by $M_B$, and write
\[
 c_B'=N-\wt(M_B)\ge0.
\]
A positive $c_B'$ is the singleton of $B$; $c_B'=0$ means that it has
none. If $A$ has no singleton, \eqref{eq:length} forces $c_B'=0$.

For every subinterval $J\subseteq I$, Lemma~\ref{lem:interval} gives
\begin{equation}\label{eq:allcounts}
 n_J(M_B)\ge n_J(D).
\end{equation}
Its polynomial identity for $I$, after division by $K_A$, is
\begin{equation}\label{eq:polynomial}
 P_D=\sum_{B\in\mathcal S}\alpha_BP_{M_B},
 \qquad \alpha_B=c_BK_B/K_A>0.
\end{equation}
We show that a particular coefficient on the left cannot be supplied
by any term on the right.

\subsection{No singleton}

By Lemma~\ref{lem:endpoints}, the bottoms of $M_B$ are
coordinatewise at least those of $D$, and its tops are coordinatewise
at most those of $D$. Extract the coefficient of $\delta^{\bots(D)}$
in \eqref{eq:polynomial}. It is $1$ on the left.
Lemma~\ref{lem:monomials} says that a term on the right can contribute
only if $\bots(M_B)=\bots(D)$. For such a term, equality of total
weights gives
\[
 \sum\tops(M_B)=N-\sum\bots(M_B)
                 =N-\sum\bots(D)=\sum\tops(D).
\]
The coordinatewise top inequalities must therefore all be equalities.
Reconstruction gives $M_B=D$, so $B=A$, a contradiction.

\subsection{A singleton}

Suppose $A=D\un\{(c)\}$. Nestedness gives a disjoint splitting
\[
 D=D_-\un D_0\un D_+,
\]
where
\[
 D_-=\{(a,b)\in D:a<c\},\quad
 D_0=\{(c,c)\in D\},\quad
 D_+=\{(a,b)\in D:b>c\}.
\]
Write $m_-=|D_-|$, $m_0=|D_0|$, and $m_+=|D_+|$.
Indeed, every nonrectangular pair avoids $c$, and every rectangular
pair belongs to exactly one of these classes.

Apply \eqref{eq:allcounts} to $[w,c-1]$, $[c,c]$, and $[c+1,u]$,
omitting empty intervals. Each $M_B$ has at least $m_-$, $m_0$,
and $m_+$ pairs in these three disjoint classes. Since their sum is
$m$, all three counts are exact and there are no other pairs. Thus
\begin{equation}\label{eq:split}
 M_B=M_-\un D_0\un M_+,
 \qquad |M_-|=m_-,\quad |M_+|=m_+.
\end{equation}
The notation $M_-$ and $M_+$ depends on $B$.
Applying the endpoint lemma within each of the lower and upper
intervals gives
\begin{equation}\label{eq:blockendpoints}
 \begin{aligned}
 b_i(M_-)&\ge b_i(D_-),& a_i(M_-)&\le a_i(D_-),\\
 b_i(M_+)&\ge b_i(D_+),& a_i(M_+)&\le a_i(D_+).
 \end{aligned}
\end{equation}

Cancel the common nonzero polynomial
$P_{D_0}=(\delta_c-\delta_{c+1})^{m_0}$ in
\eqref{eq:polynomial}. This leaves
\begin{equation}\label{eq:splitidentity}
 P_{D_-}P_{D_+}=\sum_{B\in\mathcal S}\alpha_BP_{M_-}P_{M_+}.
\end{equation}
Consider the monomial
\begin{equation}\label{eq:mixed}
 \theta=\delta^{\tops(D_-)+1}\delta^{\bots(D_+)}.
\end{equation}
The lower polynomials use only variables with index at most $c$;
the upper polynomials use only variables with index at least $c+1$.
Their variable sets are disjoint. By Lemma~\ref{lem:monomials}, the
coefficient of $\theta$ on the left is $(-1)^{m_-}$, and a term on the
right has this same coefficient precisely when
\begin{equation}\label{eq:survives}
 \tops(M_-)=\tops(D_-),\qquad \bots(M_+)=\bots(D_+);
\end{equation}
otherwise its coefficient is zero. Consequently at least one
$B\in\mathcal S$ satisfies \eqref{eq:survives}.

We finish by showing that such a $B$ must be $A$. Define
\[
 p=\sum\bots(M_-)-\sum\bots(D_-)\ge0,
 \qquad
 q=\sum\tops(D_+)-\sum\tops(M_+)\ge0.
\]
The nonnegativity follows from \eqref{eq:blockendpoints}. Equality
$p=0$ forces equality of the bottom multisets in the lower block;
equality $q=0$ forces equality of the top multisets in the upper block.
By \eqref{eq:survives},
\begin{equation}\label{eq:weights}
 \wt(M_-)=\wt(D_-)+p,\qquad
 \wt(M_+)=\wt(D_+)-q,\qquad
 c_B'=c-p+q.
\end{equation}

For the dominance comparisons, include $c_B'$ as a part even when it
is zero. The resulting multiset $\Pi_B$ has exactly $n$ entries and
the same total weight as $\lambda$. Dominance says that the sum of
its $k$ largest entries is at least the corresponding sum for
$\lambda$; equivalently, its $k$ smallest entries have sum at most
the corresponding sum for $\lambda$.

\smallskip\noindent\emph{The upper block forces $q=0$.}
First suppose $c_B'>c$. If $m_+=0$, then $q=0$ and
$c_B'=c-p\le c$, a contradiction. Otherwise put $r=2m_+>0$.
The $r$ parts of $M_+$ and the entry $c_B'$ are all at least $c+1$;
every other entry of $\Pi_B$ is at most $c$. Its $r$ largest entries
therefore have sum at most
\[
 \wt(M_+)+c_B'-(c+1)
 =\wt(D_+)-p-1<\wt(D_+).
\]
But the $r$ largest parts of $\lambda$ are exactly the parts of $D_+$.
This contradicts dominance. Thus $c_B'\le c$.
Now the $r$ largest entries of $\Pi_B$ are the parts of $M_+$,
whose sum is $\wt(D_+)-q$. Dominance forces $q=0$.
This conclusion also holds when $r=0$, since then both upper blocks
are empty. We have proved $c_B'=c-p\le c$.

\smallskip\noindent\emph{The lower block forces $p=0$.}
If $m_-=0$, this is immediate. Otherwise put $s=2m_->0$ and let
\[
 z=\max\bigl(\{\text{parts of }M_-\}\un\{c_B'\}\bigr).
\]
These $s+1$ entries are at most $c$, while every other entry of
$\Pi_B$ is at least $c$. Its $s$ smallest entries have sum
\[
 \wt(M_-)+c_B'-z=\wt(D_-)+c-z.
\]
The $s$ smallest parts of $\lambda$ have sum $\wt(D_-)$. Dominance
therefore implies $z\ge c$. Since every part of $M_-$ is strictly
less than $c$ and $c_B'\le c$, we obtain $c_B'=c$ and $p=0$.
Ties at $c$, including those in $D_0$, do not change the displayed sum.

All endpoint inequalities in \eqref{eq:blockendpoints} are now
equalities, using \eqref{eq:survives}. The multisets are nested, so
Lemma~\ref{lem:reconstruction} gives $M_-=D_-$ and $M_+=D_+$.
Together with \eqref{eq:split} and $c_B'=c$, this yields $B=A$,
the required contradiction. This completes the proof of
Theorem~\ref{thm:main}. \qed

\section{Further remarks}

The proof uses Schur positivity and dominance to restrict a hypothetical
decomposition, and then works entirely with Jacobi--Trudi factors.
The interval lemma requires no nestedness. Nestedness enters in the
bound of one singleton, the splitting around that singleton, and the
reconstruction of pairs from their endpoints. Dominance supplies the
length bound and the two comparisons at the end of the proof.

The last comparisons cannot simply be omitted. For example,
\[
 A=\{(5,4),(3),(2,1)\},\qquad
 B=\{(4,4),(4),(2,1)\}
\]
are nested, satisfy the interval count inequalities and the endpoint
conditions \eqref{eq:survives}, but $\phi(B)$ does not dominate
$\phi(A)$.

The argument rules out every positive decomposition and therefore
implies the existence of separating linear functionals by finite-dimensional
separation. It does not supply an explicit formula for one, or prove the
support restriction in White's Conjecture~11 on
$[\lambda,\lambda^{++}]$ \cite{White}. For fixed $e$, Taylor coefficients
of $\psi_{t,e}$ are linear functionals on $\Lambda_N$; understanding their
Schur supports may help relate the present proof to that stronger question.
The argument also gives no classification for factors with three or more
parts: the two-part factorization \eqref{eq:factor} is essential here.

\begin{thebibliography}{9}
\bibitem{White}
D.~E.~White, \emph{The Schur cone and the cone of log concavity},
\href{https://arxiv.org/abs/0903.2831v2}{arXiv:0903.2831v2} (2014).
\bibitem{GMTWYZ}
C.~Gaetz, K.~Meyer, K.~Y.~Tam, M.~Wimberley, Z.~Yao, and H.~Zhu,
\emph{Extreme rays of the $(N,k)$-Schur cone},
\href{https://arxiv.org/abs/1409.4859}{arXiv:1409.4859};
withdrawn by the authors in 2016.

\bibitem{OPAC} D.~White, {\it Open Problems in Algebraic Combinatorics}, blog maintained by S.~Hopkins, \href{https://realopacblog.wordpress.com/2019/09/22/the-schur-cone-and-the-cone-of-log-concavity/}{realopacblog.wordpress.com/2019/09/22/the-schur-cone-and-the-cone-of-log-concavity/} (2019).



\end{thebibliography}
\end{document}
